% Generated draft: requires mathematical and source review.
\documentclass[11pt]{article}
\usepackage[margin=25mm]{geometry}
\usepackage{fontspec}
\setmainfont{Latin Modern Roman}
\usepackage{amsmath,amssymb,mathtools,unicode-math}
\setmathfont{Latin Modern Math}
\usepackage{enumitem,needspace,xurl,hyperref}
\setlength{\parindent}{0pt}
\setlength{\parskip}{5pt}
\setlength{\emergencystretch}{3em}
\newcommand{\N}{\mathbb{N}}
\newcommand{\Z}{\mathbb{Z}}
\newcommand{\Q}{\mathbb{Q}}
\newcommand{\R}{\mathbb{R}}
\newcommand{\C}{\mathbb{C}}
\newcommand{\F}{\mathbb{F}}
\newcommand{\A}{\mathbb{A}}
\newcommand{\PP}{\mathbb{P}}
\newcommand{\E}{\mathbb{E}}
\newcommand{\eps}{\varepsilon}
\newcommand{\dd}{\,\mathrm{d}}
\begin{document}
\section*{UnsolvedMath ID: 30006332}
\textbf{Source problem code:} catalogue entry
\subsection{Definitions and mathematical statement}
\textbf{Unrecoverable mathematical statement.} The supplied record contains a damaged table-of-contents fragment, beginning with an unattached entry ``Conjecture'' pointing to page~1525, followed by talk titles, contributor names, page references, and a provenance sentence. It supplies no identifiable conjecture, mathematical assertion, definitions, hypotheses, domains, boundary conditions, or quantifiers. Consequently, a precise mathematical statement cannot be recovered from this evidence. None of the listed talk titles can be selected as the intended problem without additional evidence.

The supplied registry marks the statement as \texttt{unrecoverable} and provides no clean statement. Its verification note reports that the official report confirms the extraction is a damaged contents listing rather than a mathematical problem; this is a supplied verification claim, not an independent verification performed here. The registry's \texttt{open} label reflects its unresolved-claim classification rule and does not establish that any particular mathematical conjecture is open. No replacement conjecture is formulated.
\subsection{Short English statement}
This record preserves a corrupted contents listing from a 2025 Oberwolfach report. It does not identify a mathematical problem that can be stated precisely. The supplied \texttt{open} label is a curation classification; the intended claim and its mathematical status remain unidentified.
\subsection{Sources}
\begin{itemize}
\item Supplied local UnsolvedMath registry record, permanent ID \texttt{30006332}, exact source problem code \texttt{catalogue entry}, title ``Corrupted Non-Archimedean-Geometry Contents Fragment,'' in the collection ``Oberwolfach Reports Open Problems.'' The supplied proposed year is 2025; no proposer is specified. The record retains the original extraction for provenance and supplies no clean replacement statement.
\item Supplied report citation: \emph{Non-Archimedean Geometry and Applications} (2025), \emph{Oberwolfach Reports}, pp.~1517--1578, DOI \texttt{10.4171/owr/2025/29}. The source URL and the literature-source link labelled ``OWR report'' are both \url{https://doi.org/10.4171/owr/2025/29}. This authoritative report citation is supplied, but the record does not identify a recoverable problem within it.
\item Supplied dated literature curation, checked \textbf{2026-08-22}: status ``open.'' Its assessment says that the extraction is a table-of-contents fragment followed by meeting prose, identifies no mathematical claim, and is classified as open under the unresolved-claim rule. This assessment is supplied dated curation, not independently verified. The background explicitly requires independent verification of time-sensitive status and cited claims.
\item The fragment begins with ``Conjecture,'' p.~1525, without a recoverable title, author, or statement.
\item Contents attribution: Lucas Gerth, ``Moduli spaces of analytic $p$-divisible groups,'' p.~1528.
\item Contents attribution: Hélène Esnault, joint with Mark Kisin and Alexander Petrov, \emph{in progress}, ``On the vanishing in cohomology of the restriction map to the generic point,'' p.~1530.
\item Contents attribution: Ferdinand Wagner, ``$q$-Hodge filtrations \& Habiro cohomology,'' p.~1534.
\item Contents attribution: Juan Esteban Rodríguez Camargo, joint with Johannes Anschütz, Guido Bosco, Arthur-César Le Bras, and Peter Scholze, ``D-modules on arc-stacks,'' p.~1537.
\item Contents attribution: Toby Gee, joint with George Boxer, Frank Calegari, and Vincent Pilloni, ``Modularity lifting theorems for abelian surfaces,'' p.~1542.
\item Contents attribution: Thibaud van den Hove, ``Motivic Satake for ramified groups,'' p.~1544.
\item Contents attribution: Walter Gubler, joint with Joe Rabinoff, ``Psh and pluriharmonic functions on Berkovich spaces,'' p.~1547.
\item Contents attribution: Alberto Vezzani, joint with Fabrizio Andreatta and Federico Binda, \emph{in progress}, ``Relative Hyodo--Kato cohomology,'' p.~1549.
\item Contents attribution: Mattias Jonsson, joint with Sébastien Boucksom, ``Non-Archimedean methods in complex geometry,'' p.~1553.
\item Contents attribution: Gabriel Dospinescu, joint with Pierre Colmez, Wieslawa Niziol, and Juan Esteban Rodriguez Camargo, ``Completed cohomology of the Drinfeld tower for $\mathrm{GL}_{2}(\Q_p)$,'' p.~1556. The group notation is typeset from the fragment's extracted text ``GL2 (Qp )''; no mathematical assertion is inferred from this title.
\item Contents attribution: Eva Viehmann, joint with Torsten Wedhorn, ``Moduli of truncated local $G$-shtukas,'' p.~1560.
\item Contents attribution: Katharina Hübner, joint with Piotr Achinger, Marcin Lara, and Jakob Stix, ``The tame fundamental group of a rigid space,'' p.~1561. These contents attributions do not establish authorship of an intended problem.
\end{itemize}
\end{document}
